\begin{textsample}{NEG Dim 4 – 2018\_c6 – Score 12.00 – 2018\_c6/enade\_2018\_tec\_gest\_rec\_hum\_2dfa151ea0.txt}  \label{ex:f4_neg_345}
CAPÍTULO 6 ANÁLISE TÉCNICA DA PROVA Este capítulo tem por objetivo apresentar o desempenho dos estudantes concluintes de Tecnologia em Gestão de Recursos Humanos no Enade/2018.
Para isso, foram calculadas as estatísticas básicas da prova em seu todo ( seção 6.1.1 ), bem como as estatísticas relacionadas aos Componentes de Formação Geral ( seção 6.1.2 ) e de Conhecimento Específico da Área ( seção 6.1.3 ).
Dadas as suas características, foram analisadas, em separado, as questões objetivas ( seção 6.2 ) e as questões discursivas ( seção 6.3 ).
Para as questões objetivas, foram disponibilizados os índices de facilidade e de discriminação Ponto-Bisserial, também em separado, para os Componentes de Formação Geral ( seção 6.2.1 ) e de Conhecimento Específico ( seção 6.2.2 ).
De cada componente, uma das questões foi escolhida para exemplificar a análise gráfica, relacionando -se as alternativas escolhidas pelos estudantes ( inclusive o gabarito ) com o número de acertos no componente.
No Anexo I, apresenta -se a íntegra da análise gráfica para todas as questões objetivas.
Para cada uma das questões discursivas, os conteúdos dos tipos mais comuns de respostas dos estudantes são apresentados e comparados com o padrão de respostas esperado ( ver Anexo VIII com o padrão de respostas ).
Tomando -se como base as duas questões discursivas do Componente de Formação Geral, na seção 6.3.1.6 são apresentados comentários sobre a correção das respostas em relação à Língua Portuguesa.
Nas tabelas, constam as seguintes estatísticas das notas20 : média do desempenho na prova, erro padrão da média, desvio padrão, nota mínima, mediana e nota máxima para cada um de seus componentes.
Tais estatísticas contemplam o total de estudantes concluintes da Área de Tecnologia em Gestão de Recursos Humanos inscritos e presentes à prova do Enade/2018, tendo em vista agregações, ou por Grandes Regiões e o país como um todo, ou por Categoria Administrativa, Organização Acadêmica da IES e Modalidade de Ensino.
Em relação aos gráficos de distribuição de notas, o intervalo considerado foi o de 10 unidades, aberto à esquerda e fechado à direita, com exceção do primeiro intervalo, [ 0 ; 10 ], fechado em ambos os extremos.
Para os gráficos de distribuição das notas das questões discursivas, foram consideradas mais duas categorias : questão em branco21 e nota zero. 20 Uma definição dessas estatísticas pode ser encontrada no Glossário. 21 Nesse grupo estão incluídas também as respostas classificadas como nulas ou desconsideradas. 6.1 ESTATÍSTICAS BÁSICAS DA PROVA Esta seção apresenta estatísticas selecionadas e histogramas da nota geral ( 6.1.1 ) e de cada componente : Formação Geral ( 6.1.2 ) e Conhecimento Específico ( 6.1.3 ).
São também apresentadas estatísticas selecionadas de subpopulações, caracterizadas por Grande Região, Categoria Administrativa, Organização Acadêmica e Modalidade de Ensino. 6.1.1 Estatísticas Básicas Gerais Na Tabela 6.1, são apresentadas as Estatísticas Básicas da prova, por Grande Região, dos estudantes concluintes de Tecnologia em Gestão de Recursos Humanos.
A Média das notas da prova, como um todo, ( nas seções seguintes serão analisados os Componentes de Formação Geral e de Conhecimento Específico ) foi 42,2, sendo que os alunos da região Centro-Oeste obtiveram a Média mais baixa ( 40,9 ), e os da região Sul, a Média mais alta ( 42,5 ).
As demais médias foram : 41,9 nas regiões Norte e Nordeste e 42,3 na região Sudeste.
O Desvio padrão para o Brasil, como um todo, foi 12,6, sendo o maior Desvio padrão encontrado na região Centro-Oeste ( 12,9 ), e o menor, na região Norte ( 12,1 ), indicando uma dispersão um pouco menor das notas dessa última região.
A região que obteve a maior nota Máxima foi a Sudeste ( 90,7 ), ao passo que a região que atingiu a menor nota Máxima foi a Norte ( 73,8 ).
A Mediana do Brasil, como um todo, foi 42,2, sendo a maior Mediana obtida na região Sul ( 42,4 ), e a menor, na Centro-Oeste ( 41,6 ).
A nota Mínima do Brasil foi zero, obtida em quase todas as regiões, a exceção da região Norte onde foi 2,9.
Considerando -se as notas, segundo Grande Região, observa -se que existe diferença estatisticamente significativa ao nível de 95 \% apenas entre as duas maiores Médias, obtidas nas regiões Sudeste e Sul, e a menor, na região Centro-Oeste.
O comportamento das notas dos estudantes de todo o Brasil pode ser observado no Gráfico 6.1, que apresenta um histograma com a distribuição das mesmas : uma distribuição unimodal com a moda no intervalo ( 40 ; 50 ], seguido de perto pelo intervalo ( 30 ; 40 ].
Na Tabela 6.2, são apresentadas informações referentes à Média da nota final, desagregadas por Categoria Administrativa, por Organização Acadêmica e por Modalidade de Ensino.
A Média das notas da prova, como um todo, foi 42,2.
Em relação à Categoria Administrativa, os estudantes das IES Públicas obtiveram Média mais alta ( 47,2 ) que a Média nacional, e os das IES Privadas obtiveram Média mais baixa ( 42,1 ).
Observa -se que existe diferença estatisticamente significativa entre as médias das notas das IES Públicas e as das Privadas.
A diferença entre as médias das regiões Sul e Centro-Oeste ( 1,6 ), a maior e a menor Média, é inferior à diferença entre a das IES Públicas e a das Privadas ( 5,1 ), caracterizando -se maior diversidade administrativa do que regional.
O Desvio padrão para as IES Públicas ( 13,2 ) foi superior ao do Brasil, como um todo ( 12,6 ), indicando uma dispersão um pouco maior das notas nessa Categoria Administrativa.
No tocante à Organização Acadêmica, os CEFET/IFET e os Centros Universitários obtiveram Média mais alta que a nacional ( 50,1 e 42,8, respectivamente ).
A Média das Universidades ( 41,6 ) foi menor e a das Faculdades ( 42,2 ) igual à nacional.
Constata -se que existe diferença estatisticamente significativa ao nível de 95 \% entre a Média mais alta, 50,1 dos CEFET/IFET, e a dos demais tipos de Organização Acadêmica.
Também constata -se que não existe diferença estatisticamente significativa ao nível de 95 \% entre as médias mais baixas, a das Universidades ( 41,6 ) e a das Faculdades ( 42,2 ).
A Média da Modalidade Educação a Distância ( 41,4 ) foi inferior à Média da Educação Presencial ( 42,6 ), e a diferença entre elas é estatisticamente significativa, ao nível de 95 \%. 6.1.2 Estatísticas Básicas no Componente de Formação Geral Na Tabela 6.3, são apresentadas as Estatísticas Básicas em relação ao componente da prova que avalia a Formação Geral dos estudantes concluintes.
Os alunos de todo o Brasil obtiveram desempenho médio de 38,0.
Quanto à variabilidade, o Desvio padrão das notas dos estudantes do Brasil, como um todo, foi 15,5.
A maior Média foi obtida na região Norte ( 39,2 ), e a menor, na região Centro-Oeste ( 36,2 ).
As demais médias foram : 37,7 na região Nordeste e 38,2 nas regiões Sudeste e Sul.
Já o maior Desvio padrão foi obtido na região Centro-Oeste ( 15,7 ), e o menor, na região Nordeste ( 15,2 ).
Os demais desvios padrões foram : 15,4 na região Norte, 15,6 na região Sudeste e 15,3 na região Sul.
A maior nota no Componente de Formação Geral da prova do Enade/2018 foi 93,0, obtida por, pelo menos, um aluno na região Sudeste.
A menor nota Máxima foi obtida na região Norte ( 84,2 ).
A Mediana do Brasil, como um todo, foi 37,5, sendo a menor Mediana encontrada na região Centro-Oeste ( 35,3 ), e a maior, na região Norte ( 37,8 ).
A nota Mínima nessa parte foi zero em todas as regiões.
Considerando -se as notas, segundo Grande Região, observa -se que só existe diferença estatisticamente significativa entre as Médias das notas do Componente de Formação Geral obtidas na região Centro-Oeste em comparação com as regiões Norte, Sudeste e Sul.
As diferenças entre as demais Médias não são estatisticamente significativas.
No Gráfico 6.2, pode -se verificar a avaliação do desempenho dos estudantes do Componente de Formação Geral, a partir do histograma da distribuição das notas correspondentes.
A distribuição é unimodal, com moda em ( 30 ; 40 ], o intervalo anterior ( 20 ; 30 ] e o intervalo posterior ( 40 ; 50 ] tem frequência próxima ao modal.
O intervalo modal é imediatamente anterior ao modal da distribuição de notas da prova, como um todo ( Gráfico 6.1 ).
Nota -se, ainda que, no Gráfico 6.2, as notas apresentam maior dispersão do que as no Gráfico 6.1 ( distribuição das notas da prova ), confirmado pela comparação dos desvios padrões : 12,6 para a nota da prova, como um todo, e 15,5 para o Componente de Formação Geral.
Na Tabela 6.4, são apresentadas as informações referentes ao desempenho dos concluintes do Componente de Formação Geral, em diferentes agregações : Categoria Administrativa, Organização Acadêmica e Modalidade de Ensino.
Observa -se que existe diferença estatisticamente significativa entre as médias dos tipos de Categoria Administrativa : a maior Média obtida por estudantes de IES Públicas ( 45,9 ), e a menor, pelos de IES Privadas ( 37,9 ).
Nota -se que também há diferença estatisticamente significativa entre a maior das médias dos quatro tipos de Organização Acadêmica e a menor delas, sendo o valor maior para os CEFET/IFET ( 50,4 ) e a menor para as Faculdades ( 37,6 ).
No que se refere à Modalidade de Ensino, a Educação Presencial apresentou menor Média do que a da Educação a Distância ( médias 37,9 e 38,4, respectivamente ).
Tais médias não apresentam, entretanto, diferença estatisticamente significativa. 6.1.3 - Estatísticas Básicas do Componente de Conhecimento Específico Na Tabela 6.5, são apresentadas as Estatísticas Básicas referentes ao Componente de Conhecimento Específico da área de Tecnologia em Gestão de Recursos Humanos.
A Média do desempenho dos alunos do Brasil, como um todo, foi 43,6.
A maior Média foi obtida na região Sul ( 43,9 ), e a menor, na região Centro-Oeste ( 42,5 ).
As demais médias foram : 42,8 na região Norte, 43,3 na região Nordeste e 43,7 na região Sudeste.
Quanto à variabilidade das notas, o Desvio padrão do Brasil, como um todo, foi 13,9, sendo o maior Desvio padrão observado na região Centro-Oeste ( 14,2 ), e o menor, na região Norte ( 13,4 ).
Os demais desvios foram : 13,8 na região Nordeste e 13,9 nas regiões Sudeste e Sul.
A Mediana das notas dos estudantes de todo o Brasil foi 44,0, o mesmo valor da Mediana das regiões Sudeste e Sul.
A maior Mediana ocorreu nas regiões Sudeste e Sul ( 44,0, em ambas ), e a menor, na região Centro-Oeste ( 42,8 ).
As demais medianas foram : 43,2 na região Norte e 43,7 na região Nordeste.
A nota Máxima do Brasil, como um todo, foi 94,3, sendo obtida por, pelo menos, um aluno na região Sudeste.
As demais notas máximas foram : 76,3 na região Norte, 82,3 na região Nordeste, 85,3 na região Sul e 83,7 na região Centro-Oeste.
A nota Mínima foi zero em quase todas as regiões, com exceção da região Norte, que cuja nota Mínima foi 2,0.
Observa -se que existe diferença estatisticamente significativa apenas entre as Médias das regiões Sudeste e Sul com a da região Centro-Oeste, das notas do Componente de Conhecimento Específico.
Assim como nos Gráficos 6.1 e 6.2, no Gráfico 6.3, apresentado a seguir, é possível ser feita uma avaliação do desempenho de concluintes de Tecnologia em Gestão de Recursos Humanos em relação ao Componente de Conhecimento Específico, com um histograma da distribuição das notas correspondentes.
Esta também é uma distribuição unimodal, e o grupo modal é o ( 40 ; 50 ], o mesmo grupo modal da prova, como um todo, e acima do grupo modal para a Formação Geral.
Cabe salientar que os intervalos contíguos ( 30 ; 40 ] e ( 50 ; 60 ] têm frequências próximas à modal, concentrando nesses três intervalos quase dois terços das notas.
Na Tabela 6.6, é apresentada uma comparação dos resultados em relação à Categoria Administrativa, à Organização Acadêmica e à Modalidade de Ensino, agora levando -se em conta o desempenho de alunos do Componente de Conhecimento Específico da prova.
No que se refere à Organização Acadêmica, a maior Média foi a dos CEFET/IFET ( 50,0 ), vindo a seguir a dos Centros Universitários ( 44,2 ), a das Faculdades ( 43,7 ) e, por fim, a das Universidades ( 42,9 ).
O maior Desvio padrão, e acima do valor para o Brasil, como um todo, foi a das Universidades ( 14,0 ).
As Faculdades obtiveram a maior nota Máxima ( 94,3 ).
Os Centros Universitários obtiveram nota Máxima 85,3 ; os CEFET/IFET, 82,0, e as Universidades nota Máxima 84,8.
As medianas foram : 43,2 nas Universidades, a menor delas ; 44,5 nos Centros Universitários ; 44,0, nas Faculdades ; e 50,5 nos CEFET/IFET, a maior.
A nota Mínima foi zero para todas as Organizações Acadêmicas, com a exceção dos CEFET/IFET, cuja nota Mínima foi 4,0.
Observa -se que existe diferença estatisticamente significativa ao nível de 95 \% no Componente de Conhecimento Específico entre as notas dos CEFET/IFET e as das demais Organizações Acadêmicas.
No entanto, entre Centros Universitários e Faculdades não existe diferença significativa ao nível de 95 \% nesse Componente.
Quanto à Categoria Administrativa, observa -se um comportamento semelhante àquele da parte de Formação Geral e à prova, como um todo, ou seja, existe diferença estatisticamente significativa entre as médias das IES Públicas ( 47,7 ) e as das IES Privadas ( 43,5 ).
Neste caso também, a maior Média foi obtida por alunos de IES Públicas de ensino.
Quanto ao comportamento das médias no que toca à Modalidade de Ensino, há diferença estatisticamente significativa entre a Média da Educação Presencial ( 44,2 ) e a da Educação a Distância ( 42,4 ). 6.2 ANÁLISE DAS QUESTÕES OBJETIVAS Esta seção apresenta estatísticas selecionadas e histogramas das Questões Objetivas de Formação Geral ( 6.2.1 ) e de Conhecimento Específico ( 6.2.2 ).
São também apresentadas e comparadas as médias das subpopulações, caracterizadas por Grande Região. 6.2.1 Componente de Formação Geral Na Tabela 6.7, são apresentadas as Estatísticas Básicas relativas às oito questões objetivas do componente da prova que abrange a Formação Geral dos estudantes.
A Média do Brasil foi 39,6.
A menor Média foi encontrada na região Centro-Oeste ( 38,0 ), e a maior, na região Sul ( 40,4 ).
As demais médias foram : 39,7 na região Norte, 39,3 na região Nordeste e 39,6 na região Sudeste.
O Desvio padrão do Brasil foi 20,1, sendo o maior Desvio padrão encontrado na região Sudeste ( 20,3 ), e o menor, na região Nordeste ( 19,7 ).
Os demais desvios foram : 19,9 na região Norte e 20,0 nas regiões Sul e Centro-Oeste.
As medianas do Brasil, como um todo, e as das cinco regiões foram 37,5.
A nota Máxima 100,0 foi alcançada também em todas as Grandes Regiões.
As notas mínimas ( 0,0 ) foram iguais para todas as regiões.
Na Tabela 6.8, são apresentados o Índice de Facilidade e o Índice de Discriminação ( Ponto-Bisserial ) para cada uma das questões objetivas do Componente de Formação Geral.
Quanto ao Índice de Facilidade, foram usadas as seguintes cores para diferenciar o nível de dificuldade da questão : Azul para as questões classificadas com índice Muito fácil ( \&gt ; =0,86 ), verde para as questões classificadas com índice Fácil ( 0,61 a 0,85 ), amarelo para as questões classificadas com Médio ( 0,41 a 0,60 ), vermelho para as questões classificadas com Difícil ( 0,16 a 0,40 ) e roxo para as questões classificadas com Muito difícil ( \&lt ; =0,15 ).
Já quanto ao Índice de Discriminação, foram usadas as seguintes cores para qualificar a questão : As questões classificadas com índice Fraco receberam a cor vermelha ( \&lt ; =0,19 ), as classificadas com Médio receberam a cor amarela ( 0,20 a 0,29 ), as classificadas com Bom receberam a cor verde ( 0,30 a 0,39 ) e as classificadas com Muito bom ( \&gt ; =0,40 ) receberam a cor azul.
As questões objetivas do Componente de Formação Geral, segundo o Índice de Facilidade, foram assim avaliadas : das oito questões, nenhuma teve o Índice de Facilidade classificado como Muito fácil ou Fácil, por ter índice de acertos acima de 0,61.
Quatro questões foram consideradas com índice de dificuldade Médio, situando -se no intervalo entre 0,42 e 0,57 do Índice de Facilidade, ou seja, houve entre 42,0 \% e 57,0 \% de acertos, enquanto quatro questões foram classificadas na categoria Difícil, situando -se no intervalo entre 0,17 e 0,39.
Nenhuma questão apresentou menos de 15 \% de acertos, razão pela qual não houve questão classificada como Muito difícil.
O Índice de Facilidade variou de 0,17 a 0,57, e o de Discriminação, de 0,38 a 0,46.
Na Tabela 6.9, é apresentada a distribuição das questões, levando -se em conta simultaneamente a classificação dos dois índices.
As sete questões com Índice de Discriminação Muito bom figuraram entre dois níveis de dificuldade : Médio e Difícil das quais quatro na categoria Médio ( questões 2, 3, 7 e 8 ), e três, na categoria Difícil ( questões 1, 4 e 5 ).
Em particular, a questão 7 foi a que apresentou o maior poder discriminatório, com índice 0,46, e foi considerada Médio em termos de facilidade, com uma proporção de 0,43 acertos.
O máximo de acertos foi alcançado pela questão 2, com um Índice de Facilidade de 0,57.
As questões 1, 4, 5 e 6 foram consideradas Difícil, com índices de Facilidade 0,30, 0,17, 0,32 e 0,39, respectivamente.
Já os índices de Discriminação foram Muito bom para três questões ( 0,42, 0,40 e 0,44, respectivamente ) e Bom para a questão 6 ( 0,38 ).
A questão de número 4 apresentou um Índice de Facilidade de 0,17 ( Difícil ), ou seja, menos do que um quinto dos estudantes conseguiu resolvê -la, dentro do universo de participantes.
Seu Índice de Discriminação foi Muito bom ( 0,40 ).
A questão 6, também considerada Difícil ( 0,39 ) quanto ao Índice de Facilidade, obteve Índice de Discriminação Bom, 0,38.
No Gráfico 6.4, para exemplificar, analisa -se o comportamento da questão de número 4 de Formação Geral.
Trata -se de uma questão considerada Médio, em relação à facilidade e a que obteve o segundo maior Índice de Discriminação dessa parte da prova.
Neste gráfico, cada uma das cinco curvas representa o percentual de respostas em determinada alternativa da questão, em função do número de acertos dos estudantes nessa parte da prova ( Formação Geral/Múltipla Escolha ), antes de possíveis eliminações pelo critério do Ponto-Bisserial.
A curva em verde corresponde à alternativa B, a alternativa correta para essa questão.
Observa -se que entre os estudantes com menor número de acertos, nessa parte do exame, a situação mais frequente foi a escolha da alternativa C ( em preto ), incorreta.
Por exemplo, entre os estudantes que acertaram três questões, 38,3 \% escolheram a alternativa C ( em preto ), 22,4 \% escolheram a alternativa E ( em vermelho ), 17,1 \% escolheram a alternativa D ( em laranja ), 11,1 \% escolheram a alternativa A ( azul ), e 11,0 \%, a B ( em verde ).
Entre os que acertaram três respostas das questões de múltipla escolha de Formação Geral, 0,04 \% deixaram a questão em branco, e 0,02 \% marcaram mais de uma alternativa, invalidando a questão.
À medida que o número de acertos aumenta, indicando desempenho melhor nessa parte da prova, aumenta concomitantemente a proporção de estudantes que selecionaram a alternativa correta B, atingindo 100 \% para os estudantes com oito acertos.
Essa análise permite verificar como a questão discriminou os grupos de desempenho, justificando -se o alto índice obtido na questão.
Cumpre notar que não é possível inferir deste gráfico nem o índice de facilidade ( que seria uma média da proporção ponderada pela quantidade de alunos com cada uma das notas ) nem o índice de discriminação Ponto-Bisserial, por razão equivalente.
No caso extremo, no qual a grande concentração dos acertos dos alunos fosse abaixo de quatro, o índice de facilidade seria obrigatoriamente abaixo de 20 \% ( nesse exemplo ).
Caso a concentração fosse em seis acertos ou mais, o índice seria obrigatoriamente acima de 50 \%.
Os gráficos relativos às demais questões de Formação Geral constam do Anexo I. 6.2.2 Componente de Conhecimento Específico Na Tabela 6.10, são apresentadas as Estatísticas Básicas em relação às questões objetivas do Componente de Conhecimento Específico da prova de Tecnologia em Gestão de Recursos Humanos por Grande Região.
Nessa parte da prova, nenhuma questão foi anulada pela Comissão Assessora da Área.
Além disso, como será discutido ainda nesta seção, sete questões objetivas do Componente de Conhecimentos Específicos não foram usadas no cômputo das notas por terem sido descartadas pelo critério do Ponto-Bisserial.
Assim, as notas foram calculadas com base em 20 das 27 questões objetivas de Conhecimentos Específicos.
A Média do Brasil desse componente foi 46,6.
A menor Média foi observada na região Centro-Oeste ( 45,7 ), e a maior, na região Sul ( 47,1 ).
O Desvio padrão de todo o Brasil foi 15,1, sendo o menor Desvio padrão encontrado na região Norte ( 14,6 ), e o maior, na região Centro-Oeste ( 15,3 ).
A Mediana de todo o Brasil foi 45,0 e esse mesmo valor foi encontrado em todas as regiões.
A nota Máxima da prova foi 95,0, obtida nas questões objetivas do Componente de Conhecimento Específico, por, pelo menos, um aluno da região Sudeste.
A nota Mínima foi zero em todas as regiões.
Na Tabela 6.11, são apresentados os Índices de Facilidade e Discriminação ( Ponto-Bisserial ) das questões objetivas do Componente de Conhecimento Específico, para os estudantes de Tecnologia em Gestão de Recursos Humanos.
Para facilitar a diferenciação das questões, serão usadas as mesmas cores da Tabela 6.8 para as diferentes classificações dos Índices de Facilidade e de Discriminação.
Como nenhuma questão foi anulada pela CAA, as análises dizem respeito às 27 questões válidas.
A partir do Índice de Facilidade obtido, pode -se concluir que um pouco mais da metade das questões objetivas da prova foram considerados, pelo menos, Difícil : das 27 questões válidas, 14 foram classificadas como Difícil ou como Muito difícil.
Duas questões foram classificadas como Fácil, e outras 11 consideradas como Médio.
Já quanto aos índices de discriminação das questões objetivas do Componente de Conhecimento Específico da prova, obtém -se como resultado a seguinte classificação : oito das 27 questões foram consideradas pelo menos boas, enquanto uma delas teve Índice de Discriminação Muito bom.
Assim, para um pouco mais de um quarto das questões, oito em 27, os Índices de Discriminação foram Bom ou Muito bom.
Dentre as demais, 12 delas foram classificadas como Médio, e outras sete, como Fraco, sendo 19, por conseguinte, a quantidade de questões nos dois patamares mais baixos de discriminação.
Constata -se, assim, que a prova – no que se refere ao Componente de Conhecimento Específico, possuía média capacidade de discriminar entre aqueles que dominam ou não o conteúdo.
O Índice de Facilidade variou de 0,13 a 0,72, e o de Discriminação, de 0,05 a 0,42.
Na Tabela 6.12, é apresentada a distribuição das questões levando -se em conta, simultaneamente, a classificação dos dois índices.
Entre as questões que alcançaram os maiores índices de discriminação, uma delas foi classificada com Índice de Discriminação Muito bom, a de número 35.
Seu índice foi 0,42, e, quanto ao índice de Facilidade, foi classificada como Fácil, com 70 \% dos estudantes marcando a opção correta.
A questão de número 14 foi a mais difícil entre as 27 questões específicas válidas, com baixo Índice de Facilidade, apenas 13 \% de acertos.
Essa questão apresentou poder discriminatório muito baixo, 0,10, o que comprova ter sido a mais difícil para os estudantes.
Destaca -se, também, a questão 13, com Índice de Facilidade 0,24, o que, em termos percentuais, corresponde a 24 \% de estudantes que responderam, acertadamente, e o seu Índice de Discriminação foi 0,05.
Tais questões foram, portanto, pelo critério Ponto-Bisserial, consideradas inadequadas.
Por isso, as questões 13 e 14, além das questões 9, 11, 15, 22 e 29, foram eliminadas do cômputo da nota final.
A título de exemplo das análises do comportamento das questões objetivas, no Gráfico 6.5, analisa -se a questão 35 do Componente de Conhecimento Específico.
Essa questão foi considerada, pelas respostas dos estudantes avaliados na prova, como uma questão com nível de facilidade Fácil, 0,70, ou seja, 70 \% dos estudantes assinalaram acertadamente a opção C, correspondente ao gabarito.
Como já comentado, seu Índice de Discriminação foi igual a 0,42, classificado como Muito bom, o maior valor de discriminação.
Neste gráfico, cada uma das cinco curvas representa o percentual de respostas em determinada alternativa da questão 35, em função do número de acertos dos estudantes nessa parte da prova, antes de possíveis eliminações de questões pelo critério do Ponto-Bisserial.
A alternativa correta, C, representada no gráfico pela curva em preto, foi escolhida em maiores proporções pelos alunos com desempenho melhor nessa parte da prova.
Já as alternativas incorretas, também denominadas distratores, foram selecionadas, principalmente, por aqueles com notas mais baixas.
Observa -se que a soma não é 100 \%, por causa das questões não respondidas ou com mais de uma opção marcada.
Aqueles com nota zero, na sua totalidade, escolheram as alternativas D ( em laranja ) ou E ( em vermelho ).
A proporção de alunos que selecionou a resposta correta C aumenta gradativamente, chegando a atingir 100 \% para 19 acertos ou mais, enquanto a proporção dos que escolheram alternativas incorretas decai, a partir de dois ou três acertos, em função do número de acertos nessa parte da prova.
Os gráficos relativos às demais questões do Conhecimento Específico constam do Anexo I. 6.3 ANÁLISE DAS QUESTÕES DISCURSIVAS Esta seção apresenta estatísticas selecionadas e histogramas das Questões Discursivas de Formação Geral ( 6.3.1 ) e Conhecimento Específico ( 6.3.2 ).
São também apresentadas e comparadas as médias de subpopulações, caracterizadas por Grande Região. 6.3.1 Componente de Formação Geral As análises dos resultados de desempenho dos estudantes de Tecnologia em Gestão de Recursos Humanos nas duas questões discursivas relativas à Formação Geral encontram -se na Tabela 6.13 e no Gráfico 6.6.
Na Tabela 6.13, observa -se que a nota Média nesse conjunto de questões foi abaixo da obtida nas objetivas.
Os estudantes de todo Brasil obtiveram, em Formação Geral, Média 39,6 nas questões objetivas e 35,7 nas questões discursivas.
Pode -se notar também que o Desvio padrão nesse conjunto de questões foi menor do que o obtido nas objetivas : 20,1 nas questões objetivas e 18,8 nas questões discursivas.
A maior Média foi obtida na região Norte ( 38,6 ), e a menor, na região Centro-Oeste ( 33,4 ).
A Mediana de todo o Brasil, neste componente, foi 38,0, a mesma sendo obtida na região Nordeste.
Nas regiões Sul e Centro-Oeste, a Mediana foi menor ( 37,0 ), e nas demais regiões foi maior que o valor para o Brasil : Norte ( 41,5 ) e Sudeste ( 38,5 ).
A nota Máxima ( 97,0 ) foi obtida na região Sul.
Na região Norte a Máxima foi 85,5, na Nordeste, 81,0 ; e nas regiões Sudeste e Centro-Oeste foi 93,0.
A nota Mínima ( 0,0 ) foi a mesma em todas as regiões do Brasil.
No Gráfico 6.6, está representada a distribuição das notas nas questões discursivas do Componente de Formação Geral.
A moda dessa distribuição ocorre no intervalo ( 40 ; 50 ], seguida pelo intervalo ( 30 ; 40 ].
Destaca -se também o intervalo [ 0 ; 10 ], com distribuição em torno de 13 \% do total de notas, sendo que no intervalo [ 0 ; 10 ] inclui -se, além da nota zero, a frequência de alunos que deixaram esse tipo de questão em branco.
Na sequência, os resultados verificados para cada uma das questões discursivas de Formação Geral serão apresentados, estabelecendo -se relações com os conteúdos abordados em cada uma delas.
Os comentários da Banca de docentes corretores, a respeito do observado na correção das respostas dos estudantes, suas impressões e conclusões serão apresentados juntamente com a análise de cada questão.
Cumpre esclarecer que, tendo -se em vista que as questões discursivas de Formação Geral são padronizadas, ou seja, constam de todas as provas, os comentários da Banca são os mesmos para todas as carreiras acadêmicas, sendo direcionados a todos os estudantes que participaram do Enade/2018.
A seguir, serão analisados os desempenhos da Área de Tecnologia em Gestão de Recursos Humanos nas duas questões discursivas de Formação Geral do Enade/2018, comparando -se os resultados obtidos, acompanhado de comentários para cada questão. 6.3.1.1 Análise de Conteúdo da Questão Discursiva 1 do Componente de Formação Geral Os dados de Tecnologia em Gestão de Recursos Humanos, obtidos a partir das respostas à questão 1, encontram -se na Tabela 6.14 e no Gráfico 6.7.
Nessa questão, a de melhor desempenho entre as duas de Formação Geral, os alunos de todo o Brasil tiveram Média 33,2.
A maior Média para a questão 1 foi obtida na região Norte ( 38,0 ), e a menor, na região Centro-Oeste ( 30,5 ).
Quanto à variabilidade das notas, o Desvio padrão de todo o Brasil foi 23,3.
O menor Desvio padrão foi obtido na região Sul ( 22,2 ), e o maior Desvio padrão foi obtido na região Norte ( 24,3 ).
A Mediana do Brasil e de quase todas as regiões foi 30,0.
A exceção foi a região Norte, com Mediana 40,0.
As notas máxima e mínima da questão discursiva 1 foram as mesmas para todas as regiões do Brasil, 100,0 e zero.
A exceção foi a região Norte, com nota máxima igual a 90,0.
No Gráfico 6.7, mostra -se a distribuição das notas na questão discursiva 1 do Componente de Formação Geral.
Observa -se que a maior frequência corresponde ao intervalo de ( 30 ; 40 ].
A frequência desse intervalo modal é seguida bem de perto por dois outros intervalos ( 40 ; 50 ] e ( 50 ; 60 ].
Vale a pena ressaltar que um pouco mais de 12 \% dos alunos deixaram a resposta a essa questão em branco. 6.3.1.2 Comentários sobre a correção de Conteúdo das respostas à Questão Discursiva 1 A Questão 1 exigia do aluno um posicionamento sobre conteúdo humanista e crítico.
O comando da questão solicitava a redação de um texto que formulasse e articulasse argumentos consistentes envolvendo temas previstos na Portaria nº 444, de 30 de maio de 2018, que dispõe sobre o componente de Formação geral do Enade 2018, tais como : solidariedade/violência, tolerância/intolerância, inclusão/exclusão, sexualidade, relações de gênero e relações étnico-raciais ( art. 7º, item VIII ), bem como ética, democracia e cidadania ( art. 7º, item I ).
O enunciado da questão buscava o estabelecimento da relação de três dimensões – os Direitos Humanos, a Democracia e as Fake News.
Para isso, recorreu a quatro suportes : a reprodução de uma foto e três fragmentos adaptados de textos que foram disponibilizados em páginas da Internet ( da Anistia Internacional, do Jornal El País-Brasil e do Labic-Laboratório de estudos sobre Imagem e Cibercultura ).
Os fragmentos traziam dados sobre o aumento do número de assassinatos de grupos minoritários, as origens e o trabalho da vereadora Marielle Franco, sua luta pelos Direitos Humanos, seu assassinato, e a produção de fake news.
O comando solicitava um texto abordando dois tópicos : o tensionamento entre a defesa dos Direitos Humanos e a produção de notícias falsas, tendo como base a vereadora Marielle – a defesa dos Direitos Humanos realizada por ela e as notícias falsas divulgadas após a sua morte.
As fake news também estavam presentes no segundo tópico de análise : os prejuízos que provocam para a sociedade democrática.
Por um lado, os textos selecionados para constituir o enunciado da questão continham apoio suficiente para formulação e articulação de argumentos para discorrer sobre o tensionamento entre a atuação da Marielle em defesa dos Direitos Humanos e fake news.
Por outro, havia ênfase nos aspectos biográficos e na trajetória da vereadora Marielle, o que levou a respostas pouco focadas no comando, que discutiam a vida e a morte da vereadora, sua importância como um ícone dos Direitos Humanos, opiniões em torno do seu assassinato, até mesmo sobre a responsabilidade pelo crime, muitos textos com a presença/defesa de fake news ou produção de novas.
O foco sobre Marielle também ensejou respostas basicamente centradas na questão da violência.
Foram frequentes as interpretações com viés político-partidário, não previstas no padrão de resposta.
O item do comando que solicitava a discussão sobre democracia e fake news encontrou menos apoio nos textos.
Apesar de se esperar que graduandos tivessem maturidade, inclusive política, para demonstrar alguma compreensão dos aspectos principais que envolvem o convívio democrático, esse aspecto foi pouco discutido e, quando presente, quase sempre foi apenas mencionado, numa quase reprodução do comando.
A falta de entendimento do que é o conceito de Direitos Humanos também foi percebida nas respostas, seja para recomendá -lo como importante para a sociedade, seja para criticá -lo como algo que defende bandidos.
Foi marcante a ideia que vincula Direitos Humanos a direitos apenas das minorias.
Não foi incomum encontrar respostas que personificavam os Direitos Humanos, tratando -os como sujeito de alguma ação que deveria ser implementada na sociedade em geral ou em casos específicos.
As fake news foram citadas muitas vezes, mas pouco explicadas, conforme previsto no padrão de resposta.
O argumento utilizado geralmente vinculava as consequências de fake news para a imagem da Marielle, ou seja, centrava no fato em si, já indicado no enunciado, sem qualquer análise ou aprofundamento próprio.
É importante salientar que muitas respostas, provavelmente de grupos vinculados a cursos específicos, continham argumentações recorrentes : os Direitos Humanos como algo vinculado à Constituição e a tratados dos quais o Brasil é signatário ; as fake news como produto do mundo digital e globalizado, relacionado à era da pós-verdade ; a discussão sobre o papel das mídias sociais e da imprensa.
Como já apontado, houve baixa quantidade de respostas contemplando o segundo tópico do comando – prejuízos da produção de notícias falsas para a sociedade democrática ; fato motivado, possivelmente, pela falta de um elemento que ressaltasse esse ponto nos textos motivadores.
De modo geral, os estudantes não conseguiram recorrer a conhecimentos tácitos sobre democracia para construir argumentos próprios.
As respostas que abordavam o tema, geralmente estavam vinculadas à ideia do direito à verdade, do direito à integridade da pessoa ou dos direitos gerais garantidos pela Constituição.
A relação da questão com o momento político brasileiro foi indicada em um número expressivo de respostas.
Por vezes, apenas questionando a temática da questão – caso Marielle – em uma prova para avaliar o ensino superior no Brasil ou a importância dos Direitos Humanos na atual conjuntura política do país, levando à anulação da resposta.
Houve também interpretações das consequências do uso das fake news no processo eleitoral de 2018, influenciando a escolha consciente e livre das pessoas.
As respostas demonstraram atenção político social frente a esses temas, porém com pouco domínio conceitual do campo de conhecimento sobre Direitos Humanos, fake news e democracia.
Em termos da linguagem adotada pelos alunos, há dois grandes grupos de respostas.
No primeiro grupo, os que responderam de forma clara, com autonomia de pensamento e explicitação de justificativas para os aspectos solicitados no enunciado, demonstrando conhecimento sobre o que significam Direitos Humanos, fake news e democracia, além da articulação entre essas questões.
Por outro lado, um segundo grupo dos que não conseguiram ir além do que foi citado no enunciado, apenas repetindo -o, muitas vezes desviando o foco da questão para temas correlatos, como a vida de Marielle ou a violência.
Nesse segundo grupo, mais facilmente foram encontradas respostas desconexas e mal articuladas, com textos confusos que explicitam o não conhecimento sobre o tema e a falta de adequação à linguagem, necessária a um concluinte do ensino superior.
As abordagens dos Direitos Humanos mais recorrentes nas respostas corretas foram : a referência da Marielle como representante/ícone dos DH ; e os DH como direitos previstos na Constituição e em tratados internacionais dos quais o Brasil é signatário.
No entanto, também foram encontradas críticas aos Direitos Humanos, tais como : DH vistos como “ direitos de bandidos ” ; a defesa pouco embasada de frases de efeito como “ direitos humanos para humanos direitos ”.
Também chamou atenção o tratamento dos DH como um personagem, ou um órgão público, evidenciando desconhecimento sobre o tema.
Em relação às fake news, dois argumentos previstos pelo padrão de respostas foram os mais presentes : a noção de que o mundo globalizado e a facilidade de reprodução de informações pelas mídias atuais levam à criação de mentiras que são compartilhadas e que geram consequências ; a Marielle como alvo de fake news que tentaram desconstruir sua luta e seu legado.
Também foi frequente a ideia de que a produção de informações falsas é realizada por uma imprensa tendenciosa com intenção de manipular a informação.
No item relacionado a fake news e democracia, as respostas corretas afirmavam que as fake news inviabilizam o acesso à verdade e à garantia da integridade de um Estado Democrático de Direito.
Também foi abordada a influência das fake news nos resultados das eleições dos últimos anos, nos Estados Unidos e no Brasil, gerando consequências para a vida das pessoas.
Por fim, um outro aspecto abordado nas respostas a este tópico foi a polarização da sociedade em termos políticos, com a afirmação de que as fake news dificultam o debate, estimulam posições extremas e rivalidades prejudiciais à convivência em uma democracia plena.
A notas mais fracas ( de zero a 35 ) foram cerca de 20 \% das respostas corrigidas – as que não estavam em branco nem foram desconsideradas ou anuladas.
Foram respostas que se restringiam a cópia ou paráfrases do enunciado ; as que tinham foco em desdobramentos políticos e policiais do caso Marielle ; as que tinham argumentação confusa e/ou incompleta.
A grande maioria das notas, em torno de 65 \%, foi mediana, entre 40 a 70.
Nesse grupo ficaram as respostas que continham alguma abordagem adequada sobre Direitos Humanos, sobre as atividades da vereadora Marielle e sobre as fake news divulgadas após sua morte.
As notas mais altas deste grupo foram para respostas que incluíram a questão do tensionamento, quase sempre, em um texto coerente e consistente.
Muitas vezes, a nota mediana foi consequência do baixo índice de respostas à segunda parte do comando, que solicitava a discussão sobre os prejuízos das fake news para a democracia.
As melhores notas só foram obtidas por aqueles que incluíram a discussão sobre os prejuízos das fake news para a democracia.
Apenas 15 \% das respostas corrigidas receberam avaliação igual ou superior a 75.
A deficiência principal detectada pelos corretores foi a dificuldade de grande parte dos alunos para interpretar o enunciado da questão e para se expressar em textos consistentes.
De modo geral, os temas foram tratados de forma fragmentada e pouco analítica.
Além disso, chamou à atenção a predominância do lugar comum na construção argumentativa das respostas, a superficialidade do debate em torno da democracia e o desconhecimento de conceitos como Direitos Humanos, esses confundidos com instituições e até pessoas. 6.3.1.3 Análise de Conteúdo da Questão Discursiva 2 do Componente de Formação Geral A Tabela 6.15 mostra que o desempenho médio dos estudantes na questão discursiva 2 ( média 30,3 ) foi inferior ao obtido na questão discursiva 1 ( média 33,2 ).
A região Norte foi aquela cuja Média, também nessa questão, foi maior ( 32,2 ), enquanto a de menor Média foi a região Centro-Oeste ( 29,1 ).
Quanto à variabilidade das notas, o Desvio padrão de todo o Brasil foi 21,9, inferior ao obtido na questão discursiva 1 ( 23,3 ).
O maior desvio nessa questão foi obtido na região Norte ( 23,3 ), enquanto o menor foi obtido na região Sudeste ( 21,8 ).
A Mediana de todo o Brasil foi 30,0, a mesma de quase todas as regiões, a exceção da região Norte ( 35,0 ).
As notas : máxima ( 100,0 ) e mínima ( 0,0 ), foram as mesmas em todas as regiões, sem exceção.
No Gráfico 6.8, mostra -se a distribuição das notas na questão discursiva 2 do Componente de Formação Geral.
Observa -se que a maior frequência foi do intervalo ( 30 ; 40 ], com pouco mais de um quinto dos participantes ( 20,9 \% ), mesmo intervalo da questão discursiva 1.
A frequência dos alunos que deixaram a questão em branco é de um pouco mais de 12 \%, tal qual na questão discursiva 1.
Destacam -se, também, os intervalos ( 40 ; 50 ], ( 20 ; 30 ] e ( 50 ; 60 ], com frequências de 16,1 \%, 14,5 e 12,9 \%.
Nota -se, ainda, que as notas ficaram menos dispersas em comparação às da questão discursiva de número 1, o que pode ser constatado também pela comparação do desvio padrão das notas da questão discursiva 2 ( 21,9 ) e o das notas da questão discursiva 1 ( 23,3 ). 6.3.1.4 Comentários sobre a correção de Conteúdo das respostas à Questão Discursiva 2 Os textos e o comando da Questão 2 de Formação Geral são perfeitamente adequados ao que se espera avaliar em Formação Geral, ao tratar das funções sociais dos museus, contextualizando a temática sob o ponto de vista da memória e da perspectiva de futuro.
Na questão, estavam presentes, principalmente, o conteúdo humanista e crítico, a interpretação e a produção de textos sobre temas como cultura e arte, educação, ciência e democracia, previstos na Portaria que dispõe sobre o componente de Formação Geral do Enade 2018.
Os conhecimentos exigidos para responder ao comando, explicitados no padrão de resposta, vão além do senso comum, por se exigir a compreensão de conceitos como “ função social ” ( de instituições sociais ), “ memória ” ( no sentido de memória social ), “ perspectiva de futuro ” ( no contexto sócio-histórico-cultural ) e “ sociedade contemporânea ” ( que deveria independer da interpretação de cada grupo social a respeito da sociedade na qual vive ).
No entanto, apesar de envolver tais conceitos, espera -se que alunos concluintes de cursos superiores de graduação sejam capazes de discutir funções de instituições sociais, especialmente porque o enunciado trazia fato recente e de grande repercussão como apoio.
O padrão de resposta elencou as funções sociais que poderiam ser abordadas nas respostas dos estudantes, como também argumentos que poderiam ser usados ao longo das mesmas.
A tarefa dos estudantes era a construção de um texto expositivo que articulasse a importância de museus, suas funções sociais e as perspectivas de memória e de futuro, o que se mostrou de dificuldade de média a difícil.
Apesar de o enunciado ser claro, muitos estudantes, provavelmente influenciados pelos textos, limitaram -se a discutir aspectos relacionados ao Museu Nacional e, em particular, ao incêndio que o destruiu.
Foram poucos os que interpretaram corretamente o escopo do comando, que solicitava a discussão das funções sociais de museus de uma forma geral, no sentido de argumentar a importância deste tipo de instituição.
As quatro funções sociais elencadas no padrão de resposta foram encontradas na quase totalidade das redações dos estudantes de forma articulada com a importância dos museus.
Entretanto, quase nunca os textos abordavam mais do que uma função social e vinculavam, de forma clara, a importância desse tipo de instituição para a sociedade contemporânea.
Além disso, a solicitação de contextualização sob o foco museu/memória/perspectiva de futuro parece não ter sido bem compreendida e, em consequência, esteve pouco presente nas respostas.
Também se observou confusão entre os conceitos de história e memória e uma visão estreita do significado de sociedade contemporânea, às vezes, restrita à cidade do respondente, outras, ao Rio de Janeiro, ou ainda a um tipo de sociedade idealizada.
A relevância que os textos de apoio deram ao caso particular do incêndio do Museu Nacional do Rio de Janeiro e aos prejuízos causados ao seu acervo, induziu a maioria dos estudantes a elaborar respostas tratando das funções a partir da perda e, em particular, à perda do tipo de acervo daquele museu e do tipo de pesquisa que lá se desenvolve, divergindo do comando da questão, que solicitava texto abordando funções sociais de museus em geral.
Também foram encontradas muitas respostas que se limitavam a discutir o que o estudante considerava como possíveis causas do incêndio, a falta de manutenção, questões políticas e de gestão, sem responder ao que foi solicitado.
Foi também recorrente a transcrição de frases dos textos do enunciado nas respostas, diretamente ou por paráfrases, revelando certa dificuldade de criação textual.
A análise das respostas dos estudantes, além de revelar dificuldade de abordar a temática esperada, revelou outros dois aspectos predominantes.
Um deles, como já comentado, foi a ênfase no incêndio do Museu Nacional, que orientou grande parte dos textos dos estudantes.
Outro, também evidenciando desvio do esperado, foi o foco na perda patrimonial.
Nesse grupo, houve respostas que até apresentaram uma função social, mas olhada pelo ponto de vista da perda.
Nessa linha, muitas respostas enalteciam a importância do Museu Nacional, sua relevância no cenário nacional, o impacto da perda de seu acervo, isso complementado por frágeis reflexões sobre as funções sociais.
Em consequência, a perspectiva de desdobramentos mais ricos, que tratassem de aspectos sociais relacionados a diferentes tipos de museu foi prejudicada.
Muitas respostas trataram da preservação da memória e/ou guarda do acervo ; identidade, tanto individual quanto coletiva ; conhecimento sobre antepassados ; evolução da sociedade e dos seus costumes ; valorização e disseminação do conhecimento por meio da liberação do acesso aos acervos ; bem como atividades pedagógicas e exposições de arte.
O aspecto produção de pesquisa e conhecimento, embora frequente, foi mencionado de forma muito simples, sem qualquer aprofundamento, para além do que já era mencionado no enunciado.
Frequentemente apresentava -se o museu como produtor de pesquisas para o campo da saúde, o que evidencia desconhecimento das reais possibilidades de pesquisa, mesmo no Museu Nacional.
Quanto a perspectiva de futuro, quando presente no texto, era mencionada de forma breve, com afirmações genéricas sobre a possibilidade de melhoria do futuro pelo simples conhecimento do passado, numa percepção de que o passado determina o futuro, de forma linear.
Foram raras as respostas com alguma argumentação que associasse pesquisa, memória, identidade e reflexão social como base para a construção de alternativas de futuro.
A relação entre turismo e museu apareceu considerando a visita a museus como forma de lazer e, algumas vezes, como geradora de recursos.
O desempenho pode ser considerado fraco.
Muitas respostas deixavam claro que o estudante nunca tinha visitado um museu e, de uma forma geral, observou -se desconhecimento sobre suas funções sociais.
Não se esperava que história, memória social e museologia fossem assuntos dominados pelos estudantes das diversas áreas que participaram do Enade 2018, mas esperava -se que essa temática pudesse ser abordada do ponto de vista das experiências pessoais ou de conhecimento social.
Como uma vivência própria com os museus também se revelou frágil ou inexistente, o conhecimento sobre o assunto ficou bem restrito aos textos motivadores e ao senso comum, construído, basicamente, pelo noticiário em torno do incêndio.
O universo dos argumentos dos estudantes se mostrou muito limitado.
Cada dimensão foi praticamente representada por uma ideia simplista.
De forma recorrente : a preservação da memória para construção da identidade de um povo ; a compreensão do passado por meio das pesquisas como oportunidade para melhorar o futuro ; o museu como um espaço em que há um aprendizado estimulado pelo lazer.
Houve, também, uma grande concentração de respostas associando o museu a um local de guarda de acervo e de história.
As respostas à Questão 2 evidenciaram dificuldades conceituais.
Além dos conceitos de história, memória e museu se misturarem, parecendo não se distinguirem, ao mencionarem a pesquisa, também se observou falta de distinção entre pesquisas colegiais, com objetivos de aprendizado, e as realizadas por pesquisadores, com objetivos de avançar cientificamente.
A falta de compreensão do que se pedia no comando da questão e a dificuldade de expressão linguística foram dois aspectos marcantes na avaliação dos corretores.
Foi frequente os textos apresentarem ideias soltas ou de forma confusa, sem constituir unidade.
O vocabulário usado se revelou bem restrito, e fortemente apoiado no coloquial, na linguagem oral.
Predominou o senso comum em lugar de uma reflexão mais formal envolvendo os conceitos envolvidos. 6.3.1.5 Análise de Língua Portuguesa das Questões Discursivas do Componente de Formação Geral Os dados de Tecnologia em Gestão de Recursos Humanos, obtidos a partir das respostas às questões discursivas do Componente de Formação Geral, no que tange à Língua Portuguesa, encontram -se na Tabela 6.16 e no Gráfico 6.9.
Nesse aspecto, os alunos de todo o Brasil, obtiveram Média 51,7.
A maior Média em relação à Língua Portuguesa foi obtida na região Norte ( 52,9 ), e a menor, na região Centro-Oeste ( 48,1 ).
Quanto à variabilidade das notas, o Desvio padrão de todo o Brasil foi 22,4.
O menor Desvio padrão foi obtido na região Norte ( 21,8 ) e o maior Desvio padrão foi obtido na região Centro-Oeste ( 24,6 ).
A Mediana das notas de Língua Portuguesa foi 57,5 para o Brasil, como um todo, a mesma obtida nas regiões Nordeste e Sul.
A Mediana das regiões Norte e Sudeste foi 60,0, a maior entre as obtidas pelas Grandes Regiões, enquanto a da Centro-Oeste foi 55,0, a menor.
A nota Máxima para todo o Brasil foi 92,5, com, pelo menos, um aluno tirando essa nota na região Sudeste.
Nas regiões Norte e Centro-Oeste, a nota Máxima foi 85,0 ; na Nordeste, foi 87,5 ; e na Sul ; 90,0.
Já a nota Mínima foi zero em todas as regiões do país.
No Gráfico 6.9, mostra -se a distribuição das notas de Língua Portuguesa do Componente de Formação Geral.
Observa -se que a maior frequência ( quase 28 \% ) corresponde à dos alunos que obtiveram nota no intervalo ( 60 ; 70 ].
O intervalo anterior ao modal, ( 50 ; 60 ], obteve uma frequência próxima a modal, 24 \%.
Considerando os dois intervalos juntos, concentra -se mais da metade das notas.
Vale a pena ressaltar que cerca de 9 \% dos alunos que deixaram ambas as questões em branco, caracterizando -se como máximo local. 6.3.1.6 Comentários sobre a correção das respostas de Formação Geral com respeito à Língua Portuguesa As duas questões discursivas de Formação Geral ensejaram duas configurações textuais distintas, tendo em vista os seus encaminhamentos : a Questão 1 teve um caráter mais opinativo e a Questão 2, mais expositivo, conforme se explicita a seguir.
Ambas as questões permitiram a observação do desempenho linguístico dos participantes em situação formal, competência exigida no processo de formação universitária.
As habilidades de produção textual envolvidas têm, igualmente, como condição básica, o domínio do padrão formal culto da Língua Portuguesa, no que diz respeito aos aspectos relativos às questões de caráter ortográfico, textual, morfossintático e vocabular. formal culto da Língua Portuguesa, no que diz respeito aos aspectos relativos às questões de caráter ortográfico, textual, morfossintático e vocabular.
O encaminhamento das questões evidencia a expectativa de que o participante utilizasse seus conhecimentos sobre os dois temas e estruturasse seus textos de acordo com as características do registro formal adequado à situação comunicativa.
Tal configuração determina exigências quanto : à adequação da seleção vocabular, ao desenvolvimento do conteúdo, à estruturação sintática dos períodos, à organização lógica das ideias, à utilização de procedimentos de encadeamento textual e de referenciação, à obediência às exigências morfossintáticas próprias da modalidade escrita da norma-padrão, ao respeito às regras ortográficas e às regras de acentuação gráfica.
Esse encaminhamento exigiu do graduando o domínio do modo de organização textual de base dissertativo-argumentativa, essencial ao processo de formação universitária.
As habilidades de produção textual envolvidas têm como condição básica o domínio do padrão formal culto da Língua Portuguesa, no que diz respeito aos aspectos relativos às questões de caráter ortográfico, textual, morfossintático e vocabular.
O enunciado da Questão 1 tomava como base três fragmentos de textos jornalísticos publicados em sites da internet, que abordavam o aumento do assassinato de minorias e de defensores de direitos humanos no Brasil, entre eles a vereadora Marielle Franco.
O comando solicitava que o participante desenvolvesse um texto dissertativo e argumentativo com as temáticas Fake News, Democracia e Direitos Humanos.
Embora não tenha havido referência explícita à tipologia textual esperada para o desenvolvimento da resposta, essa proposta encaminhou o participante para a elaboração de um texto opinativo, com o objetivo de discutir o problema e manifestar um posicionamento crítico.
O enunciado da Questão 2 abordava o incêndio do Museu Nacional do Rio de Janeiro, ocorrido em 2018, e seus impactos na ciência brasileira e internacional, tomando como base dois fragmentos de textos jornalísticos publicados em sites da internet.
O comando solicitava que o participante desenvolvesse um texto abordando três aspectos da função social dos museus na sociedade contemporânea, encaminhando -o para a elaboração de um texto expositivo que confirmasse a posição defendida nos fragmentos motivadores a respeito da importância dessas instituições.
Para atender a essas exigências, o padrão de resposta utilizado na avaliação das questões discursivas de Formação Geral considerou aspectos relevantes ao bom desempenho linguístico como competências distintas, de modo a permitir um mapeamento detalhado do domínio dos recursos disponíveis na Língua Portuguesa para a comunicação escrita formal : ortográficos, textuais, morfossintáticos e vocabulares ( ver detalhes no Anexo VIII – Padrão de Respostas ).
Aspectos Ortográficos Para avaliar a competência relativa ao domínio das convenções ortográficas da norma-padrão da Língua Portuguesa, observou -se o respeito às regras de acentuação gráfica e da grafia padrão das palavras ( com ausência de abreviaturas próprias da linguagem da internet ), de acordo com as convenções estabelecidas pela legislação em vigor e consubstanciadas no Vocabulário Ortográfico da Língua Portuguesa, editado pela Academia Brasileira de Letras.
Aspectos Textuais Para avaliar a competência relativa ao domínio dos procedimentos de estruturação textual, procurou -se observar a utilização de mecanismos linguístico-discursivos que estabeleçam o encadeamento lógico entre as partes do texto, de forma a garantir a progressão e a coerência textuais.
Esta competência destaca -se por sua relevância, já que são avaliadas a clareza e a unidade textual.
Aspectos Morfossintáticos Para avaliar o domínio dos diferentes aspectos morfossintáticos próprios da modalidade escrita formal da norma-padrão da Língua Portuguesa, observou -se se o texto produzido atendeu às seguintes exigências : a concordância nominal, a concordância verbal, a regência nominal, a regência verbal, a flexão nominal, a flexão verbal, a correlação entre os tempos verbais, a colocação pronominal e a utilização de sinais de pontuação que contribuam para a organização lógica da frase e do texto.
Aspectos Vocabulares Para avaliar a adequação da seleção vocabular à modalidade escrita formal da Língua Portuguesa, exigida pela situação comunicativa, observou -se o atendimento às seguintes exigências : precisão na seleção/utilização do vocabulário relacionado à temática solicitada pela questão ; ausência de marcas da oralidade, como termos de sentido muito genérico ( “ coisa ”, “ negócio ”, “ você ” ) e termos de registros mais informais ( como gírias, jargões, frases feitas, ditados populares, termos regionais ).
Assim, espera -se que o participante respeite a adequação vocabular não usando gírias ou expressões coloquiais, evite repetição desnecessária de palavras e utilize um vocabulário mais formal, como solicitado por um texto dissertativo.
Para efeito de pontuação, a avaliação do desempenho linguístico considerou três grandes grupos de competências, agrupando aspectos explicitados, com os seguintes pesos relativos : aspectos ortográficos ( 20 \% ) ; aspectos textuais ( 40 \% ), aspectos morfossintáticos e vocabulares ( 40 \% ).
Os critérios de pontuação utilizados no padrão de resposta procuraram atender à diversidade de desempenho dos graduandos.
Nos textos analisados, observaram -se, assim, vários níveis de atendimento às exigências do padrão formal próprio da modalidade escrita da Língua Portuguesa, configurando um"continuum".
Por um lado, textos que traduzem com exatidão a expectativa de uso formal da língua, em todos os níveis analisados, marcados por clareza e unidade ; e, no outro extremo, textos quase incompreensíveis devido ao grande número de desvios de todos os tipos, principalmente os de caráter textual.
As duas questões discursivas de Formação Geral analisadas revelaram desempenhos distintos dos participantes.
A Questão 1 propiciou a oportunidade de maior desenvolvimento do tema solicitado, gerando consequências na estruturação textual, já que os textos foram mais longos, mais elaborados e fluentes.
Em virtude da polêmica sobre a existência das fake news e sobre a visão pública da personagem da vereadora Marielle Franco, observaram -se respostas que polemizaram o tema, favorecendo a construção argumentativa do texto.
A Questão 2 teve um comportamento distinto, devido ao caráter mais expositivo do enunciado e ao reduzido conhecimento dos participantes sobre o tema.
Assim, os textos foram mais curtos, frequentemente com repetição de ideias dos textos motivadores.
Quanto aos aspectos linguísticos analisados durante esta avaliação, observaram -se os seguintes resultados : Aspectos ortográficos : O desempenho dos participantes revelou uma diferença muito grande nos dois aspectos analisados nesta competência : baixo índice de desvios da grafia padrão e grande índice de desvios de acentuação.
Em vários casos, ocorre ausência completa de acentuação gráfica em todas as palavras do texto.
Os resultados revelaram que a tendência dominante entre os universitários brasileiros é a eliminação da acentuação gráfica, talvez motivada pelos hábitos relacionados às redes sociais e pela ausência de esclarecimento dos meios de comunicação, das autoridades e das escolas sobre as decisões do Acordo Ortográfico da Língua Portuguesa de 1990.
Os casos mais sistemáticos de eliminação do acento indicador da sílaba tônica foram : - palavras proparoxítonas ( por exemplo “ generos ”, “ arbitro ”, “ politico ”, “ publicas ”, “ numero ” ) ; - palavras proparoxítonas ( por exemplo “ generos ”, “ arbitro ”, “ politico ”, “ publicas ”, “ numero ” ) ; - palavras paroxítonas terminadas em ditongo crescente ( por exemplo “ area ”, “ noticiarios ”, “ individuo ”, “ historia ”, “ varios ”, “ noticia ”, “ negligencia ”, “ ocorrencia ” ) ; - palavras oxítonas ( por exemplo “ ninguem ”, “ esta ”, “ ate ”, “ tambem ” ).
Por outro lado, destaca -se o uso indevido do acento gráfico em palavras como “ melâncolia ” e “ intervênção ”, por exemplo.
Quanto ao domínio das convenções relativas à grafia das palavras, observaram -se desvios esporádicos como : hipercorreção pela escolha de “ e ” no lugar de “ i ”, por influência de hábitos da oralidade ( “ descriminação ” no lugar de “ discriminação ”, “ entervenção ” no lugar de “ intervenção ” ) ; desvios de grafia relacionados à variação diastrática, como por exemplo “ perca ”, “ subjulgadas ”, “ precoseito ”, “ soubre ”, “ vecendo ”, “ indesencia ”, “ apolojia ”, “ fulga ”, “ dereitos ” ;...
Observaram -se, também, muitos casos de inadequação no uso da maiúscula : uso indevido para destacar determinadas palavras-chave do texto, como “ Brasileiros ”, “ Homens ” ; ausência de maiúscula para grafar nomes próprios, como “ marielle ”.
Expressões como “ Direitos Humanos ” foram grafadas diversificadamente, com maiúsculas e com minúsculas.
Destaque -se, também, a maioria dos participantes que grafam os textos inteiramente em caixa alta, o que impede a identificação de desvios relacionados ao uso de maiúsculas.
Vale observar, também, que, ao contrário do que se esperava, não apareceram abreviaturas próprias do “ internetês ”, relacionadas ao uso de redes sociais e e-mails, como por exemplo “ vc ”.
Aspectos textuais : Para alguns dos graduandos essa competência se revelou como a mais problemática, tendo em vista os inúmeros problemas observados, desvios acumulados durante toda a formação escolar.
São eles : sequência justaposta de ideias sem encaixamentos sintáticos, provocando uma redução de estruturas subordinadas, ao lado do aumento na frequência de estruturas coordenadas e absolutas ; frases fragmentadas que comprometem a estrutura lógico-gramatical ; frases formadas apenas por oração subordinada, sem oração principal ; redução no uso de conectores para expressar relações lógicas essenciais à construção do texto como consequência da mudança de estruturação frasal ; emprego inadequado do pronome relativo ( com omissão da preposição ou a utilização do pronome relativo inadequado, como “ onde ” ), refletindo hábitos da oralidade ; ausência de recursos de referenciação, como substituição de termos por sinônimos, hiperônimos, nominalizações, expressões metafóricas.
Esses problemas revelam dificuldades em relação à estrutura formal do texto produzido, o que é preocupante ao se levar em conta que são graduandos em fase final de formação.
Em alguns textos, falta um mínimo de textualidade e de domínio do registro padrão da língua.
Na verdade, observaram -se relações linguísticas quase agramaticais, como as estabelecidas pela sequência de gerúndios sem o apoio de um ponto de partida para a organização das informações gramaticais e semânticas.
Quanto à utilização dos sinais de pontuação, observou -se uma ausência dos mesmos na maioria dos textos analisados.
Foi muito frequente a ocorrência de parágrafos sem marca interna de pontuação para separar os períodos.
São os seguintes os tipos de problemas mais encontrados : vírgula : utilização de vírgula para separar o sujeito e o predicado ; uso de vírgula no lugar do ponto para separar ideias que constituem períodos distintos ; ausência de vírgula para separar oração adjetiva explicativa ; ponto e vírgula : utilização do ponto e vírgula no lugar de vírgula ; ponto final : ausência de ponto final para separar períodos.
Aspectos morfossintáticos e vocabulares : Os resultados são transparentes em relação aos aspectos mais problemáticos no desempenho dos participantes nestas duas competências.
Em relação aos aspectos morfossintáticos, seguem algumas observações sobre os desvios mais frequentes.
Quanto à regência, o desvio mais frequente foi a falta do sinal indicativo da crase – isso revela que o usuário não tem consciência de que, sob a forma do termo “ a ”, existe a presença de uma contração entre a preposição “ a ” ( exigida pela regência do termo anterior ) e o artigo definido “ a ”.
Outro problema relacionado à regência verbal e à nominal, encontrado frequentemente nas respostas, foi a ausência de preposição antes de pronome relativo, processo generalizado na modalidade oral da língua, em situações de registro informal.
Apesar da possibilidade de que essa alteração de regência se generalize no padrão escrito da Língua Portuguesa, como já está ocorrendo até em textos jornalísticos, o não emprego da preposição foi considerado como desvio neste processo de avaliação.
Outro desvio de regência significativo foi a utilização inadequada de uma preposição ou sua ausência após o verbo ou o nome ( substantivo ou adjetivo ).
A concordância verbal e a concordância nominal apresentaram alguns desvios muito frequentes.
Quanto à concordância de número, observou -se, como apontam as pesquisas nessa área, ausência de marca ( com sujeito anteposto ou posposto ) ou uso indevido ( uso inadequado da marca de plural comandado pelo núcleo plural da locução adjetiva, apesar de o substantivo que funciona como núcleo do sintagma nominal estar no singular ).
Uma ocorrência que se destacou foi a ausência de acento circunflexo na forma plural do presente do indicativo do verbo “ ter ”, que foi considerada como um desvio na concordância verbal e não na acentuação gráfica.
Quanto à concordância de gênero, vários casos foram observados, normalmente no âmbito de sintagmas nominais longos, em que o adjetivo está afastado do substantivo.
Quanto à questão da colocação pronominal, foram poucos os desvios observados.
Concluiu -se que, no registro escrito formal, a maioria dos participantes já incorporou regras como a não introdução da frase por um pronome oblíquo e a próclise na presença de um termo atrator.
Várias marcas de oralidade foram identificadas, embora não com alta frequência : uso do pronome relativo “ onde ” como relativo universal, falta de artigo definido antes de substantivo isolado ou antes de termos em enumerações.
Quanto aos aspectos vocabulares, alguns tipos de inadequação foram observados : expressões da oralidade ; seleção vocabular incompatível com o contexto, gerando falta de inteligibilidade ; falta de domínio de vocabulário mais abstrato e de maior complexidade, essencial ao desenvolvimento do texto de base dissertativa.
O principal aspecto observado foi a excessiva repetição de certas palavras, como o termo"pessoa", por exemplo, revelando limitação de repertório vocabular.
Em função do tema solicitado na Questão 1, diferentes grafias da expressão fake news não foram consideradas.
Para sintetizar, os problemas que mais se destacaram nos textos analisados foram : a ) truncamentos sintáticos, caracterizados por estruturas frasais incompreensíveis devido à ruptura da complexidade sintática própria da modalidade escrita ; b ) redução de estruturas subordinadas, compensada pelo aumento na frequência de estruturas coordenadas e absolutas, por um lado, ou pela elaboração de estruturas textuais prejudicadas pelo excesso de ideias sem a devida conexão subordinativa ; c ) redução no uso de conectores para expressar relações lógicas essenciais à construção do texto, substituídos pela exigência de inferência por parte do interlocutor para suprir a sua ausência ; d ) redução no uso do subjuntivo, ao lado da ampliação do uso do indicativo, combinado a estruturas frasais coordenadas ou absolutas ; d ) empobrecimento do processo de referenciação, com a repetição de palavras, sem processos mais sofisticados de substituição ; e ) simplificação da marcação da categoria tempo na morfologia verbal ; f ) desrespeito às regras de concordância verbal e nominal próprias do padrão formal ; g ) ausência do acento grave indicador da crase, revelando falta de conhecimento dos critérios de regência verbal e nominal ; h ) falta de domínio de vocabulário mais abstrato e de maior complexidade, essencial ao desenvolvimento do processo dissertativo ; i ) redução do emprego da acentuação gráfica.
Como resultado da correção, merece atenção a constatação de que a modalidade escrita tem apresentado a tendência a uma evidente simplificação, aproximando -se das características da modalidade oral da Língua Portuguesa.
No caso dos textos de base dissertativa ( expositiva ) e de base argumentativa, inscritos em um registro formal, a distância entre as duas modalidades é ainda maior, o que provoca desvios recorrentes em todos os quatro aspectos analisados : ortográficos, textuais, morfossintáticos e vocabulares.
Destaca -se, nessa avaliação, o comprometimento textual, por meio de estruturas fragmentadas e/ou truncadas, rompendo a complexidade sintática esperada no padrão formal.
No Gráfico 6.10, representa -se a distribuição das notas nas questões discursivas no Componente de Conhecimento Específico.
A moda dessa distribuição ocorre no intervalo [ 0 ; 10 ], com quase 22 \% do total de participantes.
Observa -se que nesses intervalos estão computadas as ocorrências de respostas em branco para as três questões discursivas de Conhecimentos Específicos.
Vale a pena salientar que o intervalo ( 20 ; 30 ] tem uma frequência muito próxima ao modal ( 20,8 \% ).
A análise de cada uma destas questões será feita a seguir. 6.3.2.1 Análise da Questão Discursiva 3 do Componente de Conhecimento Específico Na questão 3, cujos resultados encontram -se descritos na Tabela 6.18, a Média dos estudantes de todo o Brasil foi 24,1.
A menor Média nessa questão foi obtida pelos alunos da região Sul ( 20,8 ), enquanto a maior Média foi obtida na região Norte ( 25,8 ).
Quanto à variabilidade das notas, o Desvio padrão de todo o Brasil foi 21,2.
O maior Desvio padrão foi obtido na região Nordeste ( 21,8 ), enquanto o menor foi obtido na região Centro-Oeste ( 20,5 ).
A nota Máxima, 100,0 pontos, foi alcançada por, pelo menos, um aluno das regiões Sudeste e Sul.
Nas regiões Norte e Centro-Oeste, a nota Máxima alcançada foi 80,0 e, na região Nordeste foi 90,0.
A Mediana do Brasil, como um todo, foi 20,0, a mesma para a maioria das regiões, com exceção das regiões Norte e Sudeste onde a Mediana foi maior, 30,0.
A nota Mínima ( 0,0 ) também foi a mesma em todas as regiões do Brasil.
No Gráfico 6.11, mostra -se a distribuição das notas na questão discursiva 3, do Componente de Conhecimento Específico da área de Tecnologia em Gestão de Recursos Humanos.
Sem se levar em consideração que 10 \% dos estudantes deixaram a questão em branco e que 20 \% receberam nota zero, a moda dessa distribuição seria o intervalo ( 30 ; 40 ], seguido do intervalo ( 20 ; 30 ]. 6.3.2.2 - Comentários sobre as respostas à Questão Discursiva 3 Questão 3 abordava um tema de grande importância em organizações de qualquer segmento e tamanho, que é Liderança.
Essa competência, que pode ser definida como a capacidade de influenciar pessoas no alcance de objetivos comuns, representa um desafio para as organizações atuais, pois identificar líderes significa atrair e reter, em seu quadro, pessoas com proatividade, capacidade de conduzir e desenvolver equipes que trabalham para atingir seus objetivos pessoais e os da organização, capazes, ainda, de facilitar novas conexões entre as pessoas e o relacionamento interpessoal entre a equipe e com clientes internos e externos.
Liderança Participativa e Empowerment são competências requeridas não só do gestor de Recursos Humanos, mas de gestores de todas as áreas.
A questão foi preparada de modo a possibilitar análise integrada destes conceitos ( Empowerment e Administração Participativa ) com papel do gestor e processo de fonte poder e autoridade.
O item ‘ a ’ envolvia a importância do papel do gestor, o que remete a estilos de liderança e seu papel da condução de equipes.
Esperava -se que os alunos associassem a implantação da Administração Participativa à Cultura Organizacional adequada, acompanhada de mudança de atitudes de gestores em relação a formas de reconhecimento, recompensa, comunicação ampla e aberta de objetivos e informações organizacionais, tolerância a erros, entre outras ações que promovem nos funcionários a vontade de assumir novos desafios e o sentimento de Empowerment.
O item ‘ b ’, de certa forma, exigia conhecimento mais aprofundado, por solicitar a relação entre o modelo participativo e sua articulação com os processos de fontes de poder e autoridade.
Faz, ainda, alusão aos quadrinhos que apresentam uma charge onde o gestor solicita aos funcionários que se sintam empoderados, porém, sem sucesso, pois, na ocorrência da primeira “ opinião divergente ”, o gestor opta por voltar atrás e continuar com o modelo antigo, centralizador e autoritário.
O início do comando do item ‘ b ’ ( “ Em relação ao cenário que se apresenta nos quadrinhos ” ) parece ter direcionado os alunos para a discussão sobre o tema gestão participativa e empowerment como conceitos contrários a poder e autoridade.
Talvez por isso, não foram encontradas respostas mais completas, ou que fizessem alusão aos tipos de poder ( tipologia de poder – Raven, 1965 ).
Baseando -se no fato de que era necessário conhecer e relacionar o modelo participativo ao uso das fontes de poder e autoridade, o que não é possível somente com conhecimentos superficiais sobre o tema liderança, pode -se classificar a Questão 3 como difícil para formandos dos cursos de Tecnologia em Gestão de Recursos Humanos.
O item ‘ a ’ foi respondido pela maioria dos formandos e, por isso, foi o responsável, em grande parte, pela nota média obtida, ainda que baixa ( 24,1 ).
Grande parte dos alunos mencionou o papel do gestor na aceitação de ideias, no fornecimento de feedback, na aceitação de erros, no compartilhamento das decisões, gerando desenvolvimento e motivação nos funcionários.
Nas respostas, observou -se o uso de termos específicos de forma adequada.
Poucos associaram a implantação do Empowerment à necessária mudança de cultura organizacional.
A correção do item ‘ b ’ deixou clara a maior dificuldade dos alunos em dissertar sobre o poder como atributo das pessoas no modelo participativo.
Grande parte das respostas encontradas permaneceu no curso das ideias do item ‘ a ’, abordando o papel do gestor no empoderamento dos funcionários, o modelo participativo contrário ao autoritarismo ( talvez por incentivo dos quadrinhos e da forma como a pergunta foi formulada, fazendo alusão a eles ), a necessária aceitação de erros, melhor comunicação, etc.
Pode -se dizer que a resposta esperada, explicando como são articuladas as fontes de poder e autoridade no modelo participativo, praticamente não foi encontrada.
As abordagens mais frequentes foram a associação do Empowerment à Gestão Participativa, a indicação do conceito de Empowerment como contrário ao autoritarismo e a importância da Gestão Participativa no desenvolvimento do colaborador.
Foi também citada, várias vezes, a importância da aceitação de erros, do feedback, da comunicação ampla na gestão participativa.
Já no item ‘ b ’, grande parte dos estudantes se apegou a discutir a situação apresentada nos quadrinhos para responder, esquecendo os conceitos mais densos solicitados – fontes de poder e autoridade – como esperado pelo padrão de resposta, que envolvia a menção aos tipos de poder.
O desempenho dos alunos nesta questão ficou aquém do esperado. 10 \% dos estudantes deixaram a resposta em branco.
Em torno de 73 \% do total de participantes receberam notas de zero a 40.
Essas respostas, consideradas fracas, variaram desde aquelas respondidas de forma contrária ao esperado ( nota zero ), citando, por exemplo a necessidade de rigidez da chefia ou que o chefe deve tomar todas as decisões, comandando os funcionários para não acontecer o que aconteceu nos quadrinhos, aquelas que somente descreveram o que aconteceu nos quadrinhos, até aquelas respondidas parcialmente, citando que o gestor é o responsável por dar segurança à equipe, pela implantação do modelo participativo, por facilitar a comunicação, por aceitar erros e fornecer feedback.
Foram consideradas medianas ( com notas de 45 a 70 ) cerca de 16 \% das respostas dos participantes.
Nessas os estudantes discorreram, de forma mais completa, sobre o papel do gestor na implantação do Empowerment, que é o de dar liberdade para tomar decisões e promovendo o desenvolvimento da equipe.
Além disso, abordaram a importância da comunicação ampla, a gestão participativa extinguindo a figura do chefe e colocando ênfase na figura do líder, a valorização das pessoas e o compartilhamento das decisões.
Apenas 137 ( 0,7 \% ) respostas foram consideradas boas ( com notas de 75 a 100 ).
Essas foram as respostas que mais se aproximaram do padrão de resposta esperado nos itens ‘ a ’.
No item ‘ b ’, os estudantes dissertaram sobre as fontes de poder e autoridade substituídas por modelos mais participativos dos funcionários, no conceito de Empowerment ligado ao Modelo Participativo e contrário ao autoritarismo e sobre a necessária habilidade do gestor para implantação da Gestão Participativa.
A Questão 3 pode ser considerada difícil para os estudantes, especialmente no item ‘ b ’, já que nas respostas os estudantes não demonstraram conhecimentos sobre as fontes de poder e autoridade.
Pode -se dizer que quase a totalidade dos alunos se restringiu a falar do Empowerment e da importância da Gestão Participativa.
Conceitos como o do papel do líder de influenciar e de conduzir as pessoas no ambiente organizacional, contribuindo para os relacionamentos intra e interpessoais e para o trabalho em equipe, bem como o da percepção da liderança como atributo das pessoas e não do cargo, são conhecimentos de grande importância na formação do Tecnólogo em Gestão de Recursos Humanos.
No entanto, tais conceitos não tiveram destaque nas respostas corrigidas.
De maneira geral, as respostas foram superficiais, incompletas ou incorretas. 6.3.2.3 - Análise da Questão Discursiva 4 do Componente de Conhecimento Específico Na Tabela 6.19, constam as informações relativas à questão 4 do conjunto de questões do Componente de Conhecimento Específico.
O desempenho dos estudantes nessa questão foi superior ao das questões 3 e 5.
A Média geral do Brasil foi 33,2, sendo a maior Média registrada na região Sul ( 36,7 ), e a menor, na região Norte ( 29,9 ).
A nota Máxima ( 100,0 ) foi atingida por, pelo menos, um concluinte em todas as regiões.
A Mediana em todo o Brasil foi 35,0, mesmo valor obtido nas regiões Nordeste e Sudeste.
Na região Norte a Mediana 25,0, a Sul foi 40,0 e na Centro-Oeste foi 30,0.
A nota Mínima foi zero em todas as regiões.
No Gráfico 6.12, está representada a distribuição de notas da questão discursiva 4, do Componente de Conhecimento Específico.
A moda dessa distribuição está no intervalo ( 40 ; 50 ], seguido pelos intervalos ( 30 ; 40 ] e ( 20 ; 30 ].
Vale a pena ressaltar que o intervalo modal tem frequência de 16,2 \% e a classe de estudantes que deixaram a questão em branco tem frequência de 16,0 \%.
A soma das frequências dos que deixaram a questão em branco e a dos que obtiveram nota zero atinge quase um quarto dos participantes. 6.3.2.4 Comentários sobre as respostas à Questão Discursiva 4 A Questão 4 abordava um tema de grande valor na sociedade atual e que, de forma geral, é o foco das organizações comprometidas com as pessoas : a qualidade de vida no trabalho.
Esse tema está diretamente relacionado às competências requeridas de um Administrador de Recursos Humanos : ético, humanista, que atua de forma holística no ambiente de trabalho, através de ações que promovem saúde, bem-estar e qualidade de vida.
Cabe ao gestor criar políticas organizacionais capazes de promover esse ambiente de qualidade, que interfere diretamente na satisfação das pessoas e, por conseguinte, na produtividade da organização.
Na primeira parte desta questão ( item ‘ a ’ ), era esperado que os estudantes associassem a qualidade de vida no trabalho à Teoria dos Dois Fatores de Herzberg, descrevendo três fatores dessa teoria que influenciam a promoção desse ambiente de satisfação e bem-estar.
Esse item pode ser considerado difícil, pois era necessário demonstrar conhecimento sobre a teoria e sua aplicação no ambiente organizacional.
No item ‘ b ’ solicitava -se que fossem citadas quatro ações de saúde e segurança no trabalho que contribuíssem para a melhoria da qualidade de vida no ambiente organizacional, o que, de certa forma, é do conhecimento de todas as pessoas.
Portanto, um item com baixo nível de dificuldade.
Assim, acredita -se que o nível de profundidade requerido na Questão 4 ficou equilibrado e esta pode ser considerada de dificuldade média.
O texto introdutório presente no enunciado da questão deixava clara a relação entre qualidade de vida no trabalho e os indicadores de produtividade e satisfação interna, percebidos através do número de afastamentos, acidentes e da rotatividade, facilitando, de certa forma, a elaboração de uma resposta ao item ‘ b ’.
No item ‘ a ’, foram consideradas corretas as respostas que discorreram sobre os Fatores Higiênicos ( ou Fatores de Manutenção ), que se localizam no ambiente, ou seja, fora do indivíduo, ou sobre os Fatores Motivacionais, que se relacionam aos sentimentos de realização pessoal e profissional ( reconhecimento, responsabilidade, desenvolvimento, inovação, participação, autonomia ).
No entanto, um número significativo de respostas apenas citava fatores que influenciam a qualidade de vida no ambiente de trabalho com base no senso comum, sem nenhuma relação com a teoria proposta.
No item ‘ b ’, que solicitava a menção a quatro ações de saúde e segurança no trabalho que contribuem para melhoria da qualidade de vida no ambiente organizacional, observou -se maior número de respostas dentro do esperado, pela facilidade de apontar ações corretas usando o senso comum e a vivência nas organizações.
Foi observado que, talvez devido à maior necessidade de conhecimento para elaboração da resposta do item ‘ a ’, muitos alunos se restringiram a responder ao item ‘ b ’.
Assim, boa parte das respostas obteve pontuação apenas no item ‘ b ’.
Foram, ainda, encontradas respostas em que os estudantes se basearam em fatores negativos para descrever a influência desses na qualidade de vida, tais como conflitos, realização de tarefas de outras pessoas, insatisfação com remuneração, má alimentação, tabagismo e sedentarismo.
As respostas apresentadas dessa forma não foram consideradas como corretas.
Nessa questão, 16 \% dos participantes não responderam, e foram consideradas respostas fracas ( com notas de zero a 40 ), desde aquelas em que os alunos somente reproduziram parte do enunciado, até aquelas em que apenas citavam fatores do ambiente que promovem qualidade de vida nesse ambiente de trabalho, porém, sem qualquer fundamentação baseada na teoria. 46 \% das respostas obtiveram notas nesse intervalo.
As respostas medianas ( com nota superior a 40 e até 70 ) corresponderam a 33 \% dos participantes.
Nessas, em sua maioria, os estudantes apresentavam a associação de um ou dois fatores que influenciam a qualidade de vida no ambiente de trabalho relacionadas à Teoria de Herzberg no item ‘ a ’.
E, no item ‘ b ’, citavam corretamente ações de saúde e segurança no trabalho que contribuem para a melhoria da qualidade de vida no ambiente organizacional. segurança no trabalho que contribuem para a melhoria da qualidade de vida no ambiente organizacional.
Apenas cerca de 900 respostas ( 5,2 \% do total de estudantes ) foram consideradas boas, com notas de 75 a 100.
Essas foram aquelas em que os estudantes responderam ao esperado no item ‘ a ’, descrevendo, à luz da Teoria dos Dois Fatores de Herzberg, fatores que influenciam a qualidade de vida dos colaboradores no ambiente de trabalho e citando, corretamente, ações de saúde e segurança no trabalho que contribuem para a melhoria da qualidade de vida no ambiente organizacional.
Devido ao número pequeno de estudantes que respondeu corretamente ao item ‘ a ’, fazendo a associação de fatores que influenciam a qualidade de vida no ambiente de trabalho, segundo a Teoria dos Dois Fatores de Herzberg, pode -se afirmar que o conhecimento dessa teoria não foi demonstrado pela maioria dos respondentes.
Como já comentado, os estudantes tiveram maior facilidade em responder ao item ‘ b ’, pois ações que contribuem para a melhoria da qualidade de vida no ambiente de trabalho representam um conhecimento que praticamente independe da formação acadêmica, sendo conhecimento associado à vida profissional e bastante discutido pela mídia em geral. 6.3.2.5 Análise da Questão Discursiva 5 do Componente de Conhecimento Específico Na Tabela 6.20, constam as informações relativas à questão 5 do conjunto do Componente de Conhecimento Específico.
O desempenho dos estudantes de todo o Brasil nesta questão foi o mais baixo entre as três questões discursivas desse componente.
A nota Média dos estudantes de todo o Brasil foi 22,0.
A maior Média foi registrada na região Sudeste ( 23,0 ), enquanto a menor Média foi registrada na região Nordeste ( 19,2 ).
Quanto à variabilidade das notas, o Desvio padrão dos alunos do Brasil, como um todo, foi 23,6.
Enquanto o maior desvio foi encontrado na região Norte ( 25,6 ), o menor foi encontrado na região Nordeste ( 22,8 ).
A Mediana para o Brasil foi 15,0, sendo a mesma obtida nas regiões Norte, Sul e Centro-Oeste.
A Mediana foi abaixo da do Brasil na região Nordeste ( 10,0 ), a menor Mediana, e acima na região Sudeste ( 20,0 ).
O quantitativo elevado de notas baixas associadas a essa questão fica evidenciado no Gráfico 6.13.
Para o conjunto de alunos de Tecnologia em Gestão de Recursos Humanos do Brasil, as notas Mínima e Máxima foram as mesmas em todas as regiões, zero e 100,0, respectivamente.
No Gráfico 6.13, é representada a distribuição das notas da questão discursiva 5 do Componente de Conhecimento Específico.
Essa distribuição, como todas as relativas às questões discursivas dessa parte da prova, também tem modal o intervalo correspondente aos que deixaram a questão em branco, com 21,1 \% do total.
Somando -se a frequência desse intervalo com a daqueles que obtiveram nota zero, chega -se próximo a 40 \% dos participantes. 6.3.2.6 Comentários sobre as respostas à Questão Discursiva 5 A Questão 5 apresentava um contexto bastante pertinente às demandas da gestão de pessoas : o desenvolvimento humano.
Esse tema relaciona -se com as competências de inovação, dinamismo e empreendedorismo esperadas do gestor de Recursos Humanos para implementação e avaliação dos processos de Gestão de Pessoas.
Consistente e precisa, a questão foi preparada de modo a possibilitar análise integrada de aspectos organizacionais, permitindo o desenvolvimento de respostas concisas e claras, sem possibilidade de interpretação diversa.
Dentre as três questões discursivas da prova, esta pode ser considerada a mais fácil.
Apesar de o item ‘ a ’ exigir conhecimento técnico sobre as etapas do processo de treinamento e desenvolvimento, respostas corretas poderiam ser baseadas no senso comum.
No item ‘ a ’ da questão, era esperado que os respondentes descrevessem, de forma objetiva, ( e não citassem apenas ) as quatro etapas para realização de um processo de treinamento.
Já no item ‘ b ’, solicitavam -se quatro indicadores que podem ser utilizados pelo gestor de Recursos Humanos para avaliação dos resultados de um processo de treinamento e desenvolvimento de pessoas.
Considerando que no texto motivador, apresentado no enunciado da questão, relacionava -se a mudança comportamental, a aquisição e o desenvolvimento de novas habilidades e de conhecimentos para solução de problemas aos objetivos do processo de treinamento e desenvolvimento, tornava possível responder ao item ‘ b ’ usando apenas elementos daquele texto.
O maior número de respostas ao item ‘ a ’ restringiu -se a citar etapas do processo de treinamento e, em grande parte, elas não foram citadas de maneira correta.
Nas respostas a este item, vários estudantes repetiram uma mesma etapa usando palavras diferentes ( por exemplo, diagnóstico e conhecimento das necessidades ).
Outros tantos dividiam uma única etapa em duas ou mais, por exemplo, definição de quem será treinado, onde será o treinamento, número de horas, entre outras que, na verdade, são todas questões respondidas na etapa de planejamento do processo.
No item ‘ b ’, era esperado que os respondentes citassem indicadores associados a medidas de economia, de custo, tempo, qualidade na execução das tarefas e melhoria de resultados de forma geral.
Grande parte das respostas indicava mudanças promovidas pelo processo de treinamento, porém, sem estabelecer relação clara com os indicadores que gerariam melhorias.
Nessa questão, quase 60 \% dos participantes obtiveram notas de zero a 35 e 21 \% deixaram a resposta em branco.
Foram respostas em que os estudantes não demonstraram conhecimento correto de nenhuma etapa do processo de treinamento e desenvolvimento ou em que apenas citaram de uma a quatro etapas corretas no item ‘ a ’.
Quanto ao item ‘ b ’, as respostas fracas foram aquelas em que os estudantes não demonstraram conhecimento correto sobre os indicadores para avaliação de resultados do processo de treinamento e desenvolvimento de pessoas.
Cerca de 16 \% dos participantes redigiram respostas consideradas medianas, com notas de 40 a 70.
Em sua maioria, os estudantes apresentavam algumas etapas do processo de treinamento no item ‘ a ’ e, no item ‘ b ’, citavam corretamente entre um e três indicadores para avaliação do processo.
Apenas cerca 3 \% dos estudantes apresentaram respostas consideradas boas ( com notas de 75 a 100 ).
Essas foram respostas em que os estudantes responderam ao esperado nos itens ‘ a ’, descrevendo, e não apenas citando, a maioria das etapas do processo de treinamento no item ‘ a ’ e citando, corretamente, de três a quatro indicadores usados para avaliação dos resultados do processo no item ‘ b ’.
Apesar da pertinência e atualidade do conteúdo avaliado na questão, o desempenho dos estudantes evidenciou despreparo e falta de conhecimento.
Ficou clara a dificuldade de interpretação e de raciocínio lógico, além de uma frágil e imprecisa construção textual.
Tal resultado mostra um comprometimento significativo para a prática desses futuros profissionais, no que se relaciona à promoção de ações que conduzam ao desenvolvimento de pessoas no ambiente organizacional.
De forma geral, o desempenho demonstrado pelos estudantes nesta questão pode ser classificado como fraco.
Mesmo considerando as respostas incompletas ou superficiais que traziam conceitos corretos, ficou evidenciado o pouco conhecimento para descrever etapas de um processo de treinamento e para relatar indicadores que podem ser usados pelo gestor de Recursos Humanos para avaliação do processo de treinamento e desenvolvimento de pessoas. 6.3.3 Considerações Finais A prova do Enade-2018 do Curso de Tecnólogo em Gestão de Recursos Humanos, de forma geral, desafiou os alunos a pensarem e a se posicionarem como Gestores de Pessoas, principais responsáveis por ações que promovem desenvolvimento e qualidade de vida no ambiente organizacional.
Dessa forma, pode -se dizer que houve cobrança quanto à adequação dos antigos modelos de liderança a um novo ambiente, competitivo e de transformações, instalado nas atuais organizações, o que provocou rupturas e mudanças nas filosofias gerenciais.
Nesse ambiente, desponta a liderança como elemento de vital importância ao sucesso de qualquer organização, cabendo a esses líderes o desenvolvimento de competências que vão facilitar novas conexões entre pessoas, a promoção do desenvolvimento de competências técnicas e comportamentais, a melhoria da qualidade de vida no ambiente organizacional e, por conseguinte, o alcance de conquistas nos âmbitos pessoal, profissional e organizacional.
De forma geral, pode -se dizer que a correção da prova evidenciou um conhecimento deficiente, superficial e impreciso dos estudantes.
Esses demonstraram dificuldade na estruturação lógica e analítica do pensamento, além de habilidade argumentativa precária para estabelecer relações entre conceitos teórico-práticos.
O desempenho apresentado sugere que no processo de ensino e aprendizagem ainda se valoriza a memorização de conteúdos e não a construção de conhecimentos, o que pode comprometer significativamente a prática dos futuros profissionais da área.

% matched lemmas: 
\end{textsample}
